\Subhead{4.10.2}{Модуляторы и дефлекторы света, преобразователи свет--сигнал, акустооптические фильтры, анализаторы спектра и корреляторы}

В модуляторе амплитуда звука управляет $\kappa$ и дифракционной эффективностью. Выбор одного порядка превращает это в модуляцию мощности; быстрый отклик ограничен временем прохождения звука через оптическую апертуру $\tau\sim a/v_s$ и полосой преобразователя. Для бегущей волны одновременно получается оптический частотный сдвиг.

Дефлектор меняет частоту звука, следовательно $K=2\pi f_s/v_s$ и направление дифракции. В простой изотропной геометрии внешний малый угол отклонения меняется приблизительно как $\Delta\theta\simeq\lambda_0\Delta f_s/v_s$. Число разрешимых направлений порядка $N\sim\Delta f_s\tau$: широкая полоса и большая апертура увеличивают разрешение, но большая апертура замедляет переключение.

В АО-фильтре фазовый синхронизм выбирает узкую оптическую полосу; перестройка акустической частоты меняет выбранную длину волны. Анизотропное взаимодействие может одновременно менять поляризацию света. Акустическая энергия переносится по групповой скорости, поэтому снос звукового пучка влияет на реальную область перекрытия.


\paragraph{Преобразователь «свет--сигнал».}
Оптическое изображение преобразуют в последовательный электрический сигнал. В устройстве Балакшия и соавторов~\cite{balakshylightsignal}, с.~743 (описание принципа работы), короткий акустический цуг пробегает по ячейке и выделяет дифракцией узкий участок изображения. Линза собирает отклонённый свет на фотоприёмник, неотклонённый задерживается экраном. При движении цуга последовательно считываются точки строки. Время строки порядка $A/v_s$, где $A$ --- апертура; разрешение ограничено конечным размером цуга и дифракцией. Это требуемое программой направление «свет в электрический сигнал».

\paragraph{Спектроанализатор и коррелятор.}
Звуковое поле $s(t-x/v_s)$ переводит временной радиосигнал в пространственную модуляцию оптического пучка: это преобразование «сигнал–свет». Линза выполняет пространственное преобразование Фурье; положение дифрагированного пятна отображает радиочастоту, а фотодетектор измеряет спектральную интенсивность. Для коррелятора нужно оптическое перемножение с опорным сигналом/фильтром и интегрирование: результат $C(\tau)=\int s(t)r^*(t-\tau)\dd t$ имеет максимум при совпадении. Спектроанализатор интенсивности без фазовой опоры не восстанавливает произвольный комплексный временной сигнал.
